\documentclass{exam}

\newif\ifA
\newif\ifkey
\newif\ifprintselected   % required by exam_config's \ansbox

%%%%%%%%%%%%%%%%%%%%%%%%%%%%%%%%%%%%%%%%%%%%%%%%%%%%%%
%%% SET THESE IF VARIABLES
%%%%%%%%%%%%%%%%%%%%%%%%%%%%%%%%%%%%%%%%%%%%%%%%%%%%%%
% Lab Num -- Module 3's six lab periods follow L04 (Lab 1) and the two
% Module 2 measurement labs, so L23 is Lab 7. Override with \def\labNumIn{}
% if the course sequence renumbers.
\ifdefined\labNumIn
	\def\labNum{\labNumIn}
\else
	\def\labNum{7}
\fi

% is this the key or not?
\ifdefined\iskey
	\ifnum\iskey=1
		\keytrue
	\else
		\keyfalse
	\fi
\else
	%% Pick one manually
	\keyfalse % if this is NOT the key
%	\keytrue     % if this IS the key
\fi

% set up the header and footer - change for every term
\ifdefined\thisTermIn
	\def\thisterm{\thisTermIn}
\else
	\def\thisterm{Fall 2026}
\fi
%
\def\thisclass{ECE 444}

%%%%%%%%%%%%%%%%%%%%%%%%%%%%%%%%%%%%%%%%%%%%%%%%%%%%%%
%%% END SET THESE IF VARIABLES
%%%%%%%%%%%%%%%%%%%%%%%%%%%%%%%%%%%%%%%%%%%%%%%%%%%%%%
% packages and shortcuts
\input{myPackages_gr}
\usepackage[hidelinks, linkcolor=red]{hyperref}
\input{myShortcuts}
\setlength{\parindent}{0pt}
\setlength{\parskip}{6pt}

%% exam class customization
\input{exam_config}
\printselectedfalse      % full worked solutions in the key, not answers-only
\CorrectChoiceEmphasis{\color{red}}
\SolutionEmphasis{\color{red}}
\renewcommand\questionlabel{\textbf{\thequestion.}}
\renewcommand{\thepartno}{\alph{partno}}
\renewcommand{\partlabel}{(\thepartno)}

% Fillable table cell: blank in the student copy, the expected value in red on
% the KEY. Fixed-width centered columns keep the blank cells writable.
\newcolumntype{K}[1]{>{\centering\arraybackslash}p{#1}}
\newcommand{\kv}[1]{\ifkey{\color{red}#1}\fi\rule[-0.10in]{0pt}{0.34in}}

% print the answers if this is the key
\ifkey
	\printanswers % if this IS the key
\else
	\noprintanswers  %if this is NOT the key
\fi

% print the header and footer
% need to print key on top if this is the key
\ifkey
	\firstpageheader{}{\color{red}{ KEY}}{}
	\runningheader{\thisclass\hspace{1pt} -- \thisterm }{Lab \#\labNum\hspace{1pt}\color{red}{ KEY}}{\thepage\hspace{1pt} of\hspace{1pt} \numpages}
	\firstpagefooter{}{}{}
	\runningfooter{}{}{}
\else
	\firstpageheader{}{}{}
	\runningheader{\thisclass \hspace{1pt} -- \thisterm }{Lab \#\labNum\hspace{1pt}}{\thepage\hspace{1pt} of\hspace{1pt} \numpages}
	\firstpagefooter{}{}{}
	\runningfooter{}{}{}
\fi

\runningheadrule

\begin{document}
\pagestyle{headandfoot}   % exam class


%%%%%%%%%%%%%%%%%%%%%%%%%%%%%%%%%%%%%%%%%%%%%%%%%%%%%%
% title block
%%%%%%%%%%%%%%%%%%%%%%%%%%%%%%%%%%%%%%%%%%%%%%%%%%%%%%

\begin{center}
USAF ACADEMY \\
DEPARTMENT OF ELECTRICAL AND COMPUTER ENGINEERING \\
LAB \#\labNum\ (Lesson 23) - Antenna Pattern Lab \\
\textbf{Lab Sheet} \\
\thisterm \\[4mm]
\end{center}

\vspace{-4pt}
\textbf{Name:} \rule[-2pt]{2.6in}{0.4pt} \hfill
\textbf{Section:} \rule[-2pt]{0.9in}{0.4pt} \hfill
\textbf{Date:} \rule[-2pt]{1.1in}{0.4pt}

\vspace{4pt}

\fullwidth{\textbf{Learning Objective 3.6: } I can distinguish between array factor
and true antenna pattern and account for element pattern effects.}

\vspace{2pt}

This sheet is the turn-in document for the Lesson 23 lab. Report every lobe
amplitude relative to the main-lobe peak, in dBc, and state the uncertainty next
to any number read off a hand-walked trace. The lab is a team assignment, so one
sheet per team.

\begin{questions}

%%%%%%%%%%%%%%%%%%%%%%%%%%%%%%%%%%%%%%%%%%%%%%%%%%%
\question \textbf{The annotated trace.} Attach or sketch one screenshot of your
best broadside run from the Tracking tab, with the main lobe and both first
sidelobes labeled with their dBc values.

\begin{solutionorbox}[3.7in]
An acceptable trace shows one main lobe near the middle of the walk with a
sidelobe on each side of it, and all three labeled with an amplitude relative to
the main-lobe peak. The main lobe is $0$ dBc by definition, and the two
sidelobes should land within a decibel or two of $-13$ dBc. The reader has to be
able to tell which lobe each label belongs to, and to see that the values are
differences from the peak rather than absolute dBFS numbers. A trace whose
sidelobes read $-6$ or $-7$ dBc is not a uniform illumination, and it sends the
team back to Calibrate and to the element gain sliders.
\end{solutionorbox}

%%%%%%%%%%%%%%%%%%%%%%%%%%%%%%%%%%%%%%%%%%%%%%%%%%%
\newpage
\question \textbf{The amplitude comparison table.} Three lobes from the
broadside run, each as an amplitude relative to the main-lobe peak. The
mechanical column is your hand-walked trace, the electrical column is the
Rectangular-tab sweep of step 4, and the calculated column comes from theory.

\begin{center}
\small
\setlength{\tabcolsep}{4pt}
\renewcommand{\arraystretch}{1.25}
\begin{tabular}{| l | K{1.1in} | K{1.15in} | K{1.05in} | K{1.15in} |}
\hline
\textbf{Lobe} & \textbf{Mechanical (dBc)} & \textbf{Electrical sweep (dBc)} & \textbf{Calculated (dBc)} & \textbf{Mechanical $-$ electrical} \\\hline\hline
Main lobe & \kv{$0$ by definition} & \kv{$0$ by definition} & \kv{$0$} & \kv{$0$} \\\hline
First sidelobe, left & \kv{$-11$ to $-15$} & \kv{$-11$ to $-13$} & \kv{$-13.1$} & \kv{$\le 2$ dB} \\\hline
First sidelobe, right & \kv{$-11$ to $-15$} & \kv{$-11$ to $-13$} & \kv{$-13.1$} & \kv{$\le 2$ dB} \\\hline
\end{tabular}
\end{center}

\begin{solution}
The calculated $-13.1$ dBc is the $-12.8$ dB array-factor sidelobe of a uniform
eight-element array plus about $0.3$ dB of element-pattern roll-off at the
$\pm 22\mydeg$ where those sidelobes sit. A hand-walked trace in a classroom
carries about $\pm 1$ dB of multipath ripple, so the mechanical column is
accepted over a wider window than the electrical one. The two columns should
agree within about 2 dB, and a larger gap usually means the radius changed
during the walk, since received power goes as $1/r^2$ and a $10$ cm error at
$1\ \m$ is worth about $0.8$ dB. The left and right readings need not match each
other exactly; a spread of about a decibel between them is normal in a room.
\end{solution}

%%%%%%%%%%%%%%%%%%%%%%%%%%%%%%%%%%%%%%%%%%%%%%%%%%%
\newpage
\question \textbf{Written answers.} Three or four sentences each.

\begin{parts}

	\part Why do the lobe amplitudes survive a sloppy rotation while the lobe
	angles do not? Refer to what the time axis does and does not record.

	\begin{solutionorlines}[2.2in]
	The plot records received amplitude against elapsed time, and the amplitude at
	any instant is whatever the array collected from the direction the source
	occupied at that instant. Walking unevenly changes when each direction is
	visited, not how much power arrives from it, so every lobe keeps its height
	and only its position and width on the plot change. The time axis records
	nothing about where your hand was, so converting it to angle requires assuming
	a constant speed you did not have. That is why the amplitudes are reportable
	and the angles are not.
	\end{solutionorlines}

	\part Which of the four artifacts in Part 4 does an anechoic chamber remove,
	and which of them would still be present in a chamber measurement?

	\begin{solutionorlines}[2.4in]
	The absorber in an anechoic chamber removes the reflections, so the floor,
	ceiling, and whiteboard paths disappear and the $\pm 1$ dB ripple riding on
	the trace goes with them. The other three artifacts belong to the positioner
	and the range geometry rather than to the room. A hand-held source walked
	around a chamber would still have an uncalibrated angle axis, a radius that
	varies with the operator's arm, and a separation that is only as good as the
	tape measure. A range fixes those with a turntable and a shaft encoder, a
	source bolted at a fixed distance, and a documented far-field separation.
	\end{solutionorlines}

\end{parts}

\end{questions}

\vspace{1.0em}
\noindent\textbf{Documentation:} \rule[-2pt]{0.6\linewidth}{0.4pt}

\end{document}
